% Generated draft: requires mathematical and source review.
\documentclass[11pt]{article}
\usepackage[margin=25mm]{geometry}
\usepackage{fontspec}
\setmainfont{Latin Modern Roman}
\usepackage{amsmath,amssymb,mathtools,unicode-math}
\setmathfont{Latin Modern Math}
\usepackage{enumitem,needspace,xurl,hyperref}
\setlength{\parindent}{0pt}
\setlength{\parskip}{5pt}
\setlength{\emergencystretch}{3em}
\newcommand{\N}{\mathbb{N}}
\newcommand{\Z}{\mathbb{Z}}
\newcommand{\Q}{\mathbb{Q}}
\newcommand{\R}{\mathbb{R}}
\newcommand{\C}{\mathbb{C}}
\newcommand{\F}{\mathbb{F}}
\newcommand{\A}{\mathbb{A}}
\newcommand{\PP}{\mathbb{P}}
\newcommand{\E}{\mathbb{E}}
\newcommand{\eps}{\varepsilon}
\newcommand{\dd}{\,\mathrm{d}}
\begin{document}
\section*{UnsolvedMath ID: 30001565}
\textbf{Source problem code:} catalogue entry
\subsection{Definitions and mathematical statement}
A standard finite coherent configuration consists of a finite set \(X\) and a partition \(\mathcal R\) of \(X\times X\) such that the diagonal is a union of members of \(\mathcal R\), the transpose of every member belongs to \(\mathcal R\), and, for all \(R,S,T\in\mathcal R\), the number
\[
\bigl|\{z\in X:(x,z)\in R,\ (z,y)\in S\}\bigr|
\]
is independent of the choice of \((x,y)\in T\). For a coefficient field \(K\), let \(A_R\) be the adjacency matrix of \(R\), and let
\[
\mathcal A_K=\operatorname{span}_K\{A_R:R\in\mathcal R\}\subseteq\operatorname{Mat}_X(K).
\]
This is the associated matrix algebra. A representation is a unital algebra homomorphism \(\rho:\mathcal A_K\to\operatorname{End}_K(V)\). It is irreducible if \(V\ne0\) and its only invariant subspaces are \(0\) and \(V\). These definitions give a conventional framework; the supplied record does not specify the coefficient field or the representation used by the method in question.

Self-adjointness requires an adjoint operation \(*\), with \(a\) self-adjoint meaning \(a^*=a\). The record does not identify this operation. For a finite-dimensional action \(T\) of a chosen element \(a\), write \(\chi_T(t)=\det(tI-T)\). The condition that this polynomial split into distinct linear factors over a field \(L\) means
\[
\chi_T(t)=\prod_{j=1}^{\dim V}(t-\lambda_j),\qquad
\lambda_j\in L,\qquad \lambda_i\ne\lambda_j\ \text{for }i\ne j.
\]
Here \(V\) denotes the space of that action. The supplied record does not specify which action determines the characteristic polynomial or which field \(L\) is intended.

\textbf{Source question.} For representations of a coherent-configuration matrix algebra, can the desired irreducible action be computed without assuming that a chosen self-adjoint element has a characteristic polynomial splitting into distinct linear factors over the relevant splitting field?

The target irreducible action, its input description, the representation method, the permitted field extensions, and the meaning of computation are unspecified. Consequently, the record does not determine a fully quantified algorithmic problem. Its background also describes the objective as avoiding factorization of the characteristic polynomial. Avoiding a factorization step and removing the distinct-linear-factor hypothesis are different requirements; their intended relationship remains unresolved. If the relevant field is a splitting field of this very polynomial, splitting follows by definition, whereas distinctness of the roots remains an additional condition.
\subsection{Short English statement}
The question asks whether a method for computing an irreducible action of a coherent-configuration matrix algebra can dispense with a distinct-root splitting assumption on a chosen self-adjoint element. The supplied background additionally asks about avoiding polynomial factorization. The field, adjoint, input representation, target action, and relationship between these two objectives require clarification.
\subsection{Sources}
\begin{itemize}
\item Supplied local UnsolvedMath registry record, permanent ID \(30001565\), source problem code \texttt{catalogue entry}, titled ``Irreducible Coherent-Configuration Representations Without Polynomial Splitting.'' The supplied proposer field is empty; no individual proposer is identified.
\item Supplied authoritative source citation: \emph{Mini-Workshop: Exploiting Symmetry in Optimization} (2011), \emph{Oberwolfach Reports}, pp.\ 2245--2281, DOI \texttt{10.4171/owr/2010/38}. \url{https://doi.org/10.4171/owr/2010/38}. The background identifies Oberwolfach Reports as records of mathematical workshops and their open-problem discussions.
\item Supplied literature source, labeled ``original workshop report'': \url{https://doi.org/10.4171/owr/2010/38}.
\item Literature assessment checked on \texttt{2026-08-21}: the supplied curation labels the problem open and states that improving the representation method without factoring the characteristic polynomial remains open in the cited optimization literature. This is supplied dated curation, not an independently verified assessment; time-sensitive status and cited claims require independent verification.
\item The registry labels the statement \texttt{exact} and supplies a verification assertion that its mathematical content was checked against the official report, with notation, context, and question form normalized without changing the source-level claim. That assertion is recorded as supplied provenance and has not been independently verified here. The record gives proposed year 2010 and tag \texttt{owr-year:2010}, while the supplied report citation displays 2011; the date distinction remains unresolved.
\end{itemize}
\end{document}
